\section{Model and source benchmark}
\label{sec:model}

We retain the static inherited-history Hotelling environment used in the accepted manuscript of \citet{GehrigShyStenbacka2012}. The purpose is not to introduce a new switching-cost model, but to resolve the global price game generated by those primitives.

\subsection{Consumers, histories, and switching costs}

Two firms, \(A\) and \(B\), are located at the endpoints \(0\) and \(1\) of a Hotelling line. A unit mass of consumers is uniformly distributed on \(z\in[0,1]\). Each consumer purchases one unit. Gross utility is \(\beta\), transport cost is linear with coefficient \(\tau>0\), and the two firms have a common marginal cost \(c\).

Consumers arrive with inherited supplier histories. Firm \(A\) inherits consumers on \([0,x_0]\), while firm \(B\) inherits consumers on \((x_0,1]\), where
\[
x_0\in(1/2,1].
\]
Thus \(A\) is the inherited dominant firm. A consumer who changes supplier pays a switching cost \(\sigma\), and throughout the reassessment we maintain the source restriction
\[
0\le \sigma<\tau.
\]

Under uniform prices \(p_A,p_B\), an \(A\)-history consumer located at \(z\le x_0\) obtains
\[
U_A^A(z)=\beta-p_A-\tau z,
\qquad
U_B^A(z)=\beta-p_B-\tau(1-z)-\sigma,
\]
whereas a \(B\)-history consumer located at \(z>x_0\) obtains
\[
U_A^B(z)=\beta-p_A-\tau z-\sigma,
\qquad
U_B^B(z)=\beta-p_B-\tau(1-z).
\]
Full coverage is assumed, so only the identity of the chosen firm matters for demand.

It is useful to normalize net margins and switching cost by the transport scale:
\[
a=\frac{p_A-c}{\tau},\qquad
b=\frac{p_B-c}{\tau},\qquad
s=\frac{\sigma}{\tau},\qquad
d=b-a.
\]
The normalization is one-to-one for fixed \(\tau>0\), and the equilibrium prices are recovered by multiplying normalized margins by \(\tau\) and adding \(c\).

\subsection{Primitive uniform demand}

The indifferent locations for the two inherited histories are
\[
t_A(d)=\frac{d+1+s}{2},
\qquad
t_B(d)=\frac{d+1-s}{2}.
\]
The first threshold determines which \(A\)-history consumers remain with \(A\); the second determines whether any \(B\)-history consumers switch to \(A\). Because each threshold is economically relevant only inside its own inherited segment, the exact share of firm \(A\) is
\begin{equation}
m_A(d)
=
\clip\!\left(t_A(d),0,x_0\right)
+
\clip\!\left(t_B(d)-x_0,0,1-x_0\right),
\label{eq:clipped-demand}
\end{equation}
where \(\clip(y,\ell,u)=\min\{\max\{y,\ell\},u\}\). Firm \(B\)'s share is \(m_B(d)=1-m_A(d)\).

For \(0<s<1\), define
\[
L=-1-s,\qquad
\alpha=2x_0-1-s,\qquad
\gamma=2x_0-1+s,\qquad
U=1+s.
\]
Equation~\eqref{eq:clipped-demand} is equivalent to the five-piece demand correspondence
\begin{equation}
m_A(d)=
\begin{cases}
0,&d\le L,\\[0.25em]
(d+1+s)/2,&L<d<\alpha,\\[0.25em]
x_0,&\alpha\le d\le\gamma,\\[0.25em]
(d+1-s)/2,&\gamma<d<U,\\[0.25em]
1,&d\ge U.
\end{cases}
\label{eq:five-piece-demand}
\end{equation}
The middle interval is the no-switch plateau. The lower switching branch has only \(A\)-history consumers moving to \(B\); the upper branch has only \(B\)-history consumers moving to \(A\). Simultaneous switching in both directions is impossible because \(t_A-t_B=s>0\).

Firm profits are
\[
\pi_A=\tau a\,m_A(d),\qquad
\pi_B=\tau b\,m_B(d).
\]
The common factor \(\tau\) is irrelevant for best responses.

\subsection{History-based prices used for comparison}

Under history-based price discrimination (HBP), let \(p_i\) denote firm \(i\)'s price to its inherited consumers and \(q_i\) its poaching price to the rival's inherited consumers. The accepted manuscript separates weak and strong inherited dominance at
\[
\bar x(s)=\frac{3-s}{4}.
\]

On the weak-dominance branch, \(1/2<x_0<\bar x(s)\), the source price vector is
\begin{align}
\frac{p_A^d-c}{\tau}&=\frac{2x_0+1+s}{3},
&
\frac{q_A^d-c}{\tau}&=\frac{3-4x_0-s}{3},
\nonumber\\
\frac{p_B^d-c}{\tau}&=\frac{3-2x_0+s}{3},
&
\frac{q_B^d-c}{\tau}&=\frac{4x_0-1-s}{3}.
\label{eq:weak-hbp-prices}
\end{align}
These prices and their induced allocation reproduce the accepted weak-HBP benchmark.

Under strong dominance, \(x_0\ge\bar x(s)\), the source member sets \(q_A^d=c\). For the nondegenerate case \(x_0<1\), allowing prices below marginal cost creates a payoff-equivalent zero-sales family discussed in Section~\ref{sec:scope}; the nonnegative-margin convention selects the source member. The endpoint \(x_0=1\) is different: the \(B\)-history market is empty, so \(q_A\) and \(p_B\) are unused prices and the HBP price vector is not unique even under nonnegative margins. Allocation, realized profits, consumer surplus, and welfare do not depend on those unused quotes.

\subsection{Source boundary}

All equation-level comparisons in this paper use a lawful post-referee accepted manuscript of \citet{GehrigShyStenbacka2012} as the controlling source. The exact Springer typeset Version of Record has not been compared equation by equation. We therefore refer to the accepted manuscript when identifying a formula or numbered Result whose wording matters.
